\subsection{Which generalized influences could still explain full AA?}
\label{tht:generalized-influence}

Yes: an influence-based approach remains meaningful if influence
measures information revealed by queries, or changes to recoverable
latent variables. Replacing single-spin flips by bounded-size physical
updates does not suffice. These are different assertions, with different
proof obligations.

\begin{proposition}[A limitation of physically local updates]
\label{tht:local-update-obstruction}
Let \(M(x)=n^{-1}\sum_i x_i\), and let \(K(x,dy)\) be any randomized
update that changes at most \(k\) spins almost surely. Define its
unnormalized squared-change quantity by
\[
 \mathcal E_K^\mu(M)=
 \mathbb E_{X\sim\mu,\,Y\sim K(X,\cdot)}(M(X)-M(Y))^2.
\]
Then \(\mathcal E_K^\mu(M)\le4k^2/n^2\). Consequently, for the
zero-field Curie--Weiss family at fixed \(\beta>1\), no maximum of
such quantities with \(k=o(n)\) and dimension-independent normalization
can satisfy an AA lower bound in variance and polynomial degree.
\end{proposition}

\begin{proof}
The deterministic bound \(|M(x)-M(y)|\le2k/n\) implies the displayed
inequality, without any stationarity or reversibility assumption on
\(K\). In ordered Curie--Weiss, \(M\) has degree one and variance
bounded below by a positive temperature-dependent constant, whereas
the bound tends to zero. It applies uniformly to every choice of
kernel, hence also to their maximum. Multiplying by a fixed convention
factor, such as \(1/4\), changes nothing.
\end{proof}

This rules out a whole class of local Dirichlet-form repairs, including
fixed-size heat-bath block updates. It does not rule out an update that
changes a global phase while using few observed spins: physical update
size and input-query cost are distinct resources.

\textbf{Revealing information is the first viable replacement.}
For a query transcript \(\mathcal F\), use
\[
 G_\mu(Q,\mathcal F)
 =\operatorname{Var}_\mu\bigl(\mathbb E[Q\mid\mathcal F]\bigr).
\]
If \(\mathcal F_0\subseteq\cdots\subseteq\mathcal F_s\) is the
transcript of an adaptive tree and
\(m_t=\mathbb E[Q\mid\mathcal F_t]\), martingale orthogonality gives
\[
 G_\mu(Q,\mathcal F_s)=
 \sum_{t=1}^s\mathbb E(m_t-m_{t-1})^2
\]
when \(\mathcal F_0\) is trivial. This is the accumulated variance
explained by the queried coordinates.

It recognizes information missed by physical flips. For zero-field
Curie--Weiss, exchangeability and sign symmetry give exactly
\[
 \mathbb E[M\mid X_i]=v_nX_i,\qquad
 v_n=\operatorname{Var}(M),\qquad
 \operatorname{Var}(\mathbb E[M\mid X_i])=v_n^2.
\]
Thus one-spin reveal gain stays positive in the ordered phase even
though every bounded-size physical update has vanishing squared
change. For the acceptance-valued observable \(Q=(1+M)/2\), the
corresponding gain is \(v_n^2/4\).

Singleton reveal gain nevertheless cannot be the universal criterion.
Under independent fair spins,
\(Q=(1+x_1x_2)/2\) has variance \(1/4\), yet
\(\mathbb E[Q\mid X_i]=1/2\) for every individual coordinate.
Revealing \(X_1,X_2\) determines \(Q\). This is itself a one-query
quantum acceptance function: query the two-coordinate uniform
superposition with a phase oracle and measure its overlap with the
initial state. Hence even genuine low-query acceptance functions need
not have positive singleton reveal gain. A bounded-length adaptive
transcript, allowing initially uninformative queries, is the appropriate
object.

Proposition~\ref{tht:reveal-prop} already proves that polynomial
transcript length and polynomial reveal progress, uniformly under
pinnings, characterize hereditary full AA. This is a useful target,
but not a shortcut: postulating that property is essentially postulating
the desired simulation theorem. The substantive step must derive it
from independently checkable Hamiltonian structure.

\textbf{Recoverable latent-source influence is the second route.}
Suppose \(X=F(Z)\) for independent fair source bits, and each coherent
output query uses at most \(a\) source queries. Then \(Q\circ F\)
is an actual \(aT\)-query acceptance function, globally bounded and of
degree at most \(2aT\). Uniform-input AA may therefore be applied to
the source influences
\[
 I_j^{U}(Q\circ F)=\mathbb E_U\bigl[(D_j(Q\circ F))^2\bigr].
\]
A source flip may change many physical spins, so the local-update
obstruction above does not apply. For biased sources one can use the
fair-bit encodings and approximation proved in the block theorem.

Source influences alone are insufficient: the classical simulator
must operate on the actual input \(X\), not on an inaccessible sample
of \(Z\). Proposition~\ref{prop:threshold-local-posterior} supplies
the needed bridge when \(Z\mid X=x\) can be generated on demand with
few physical queries. Its approximate version also requires the
conditional mismatch bound controlling quantum hybrid error. The
whole source simulator is transported; no lower bound on a single
physical-spin influence or singleton reveal gain is asserted.

\textbf{Established versus open.} Ising forests provide an established
Hamiltonian-derived realization: independent reset/copy variables have
an edge-local posterior and a uniformly bounded conditional
continuation probability. These facts give the dimension-independent
simulation theorem, with uniform AA as its imported input. Extending
this construction, or proving transcript reveal progress directly,
for general Dobrushin interaction graphs remains open. The upstream
semigroup degree-growth and gradient-localization counterexamples
constrain possible proofs but do not refute either of these two routes.
